% =============================================================================
% Assignments/CS301/11-trees-bst-avl-coding-assignment.tex
% Coding Assignment 11: Binary Search Trees and Traversals
% (Week 11, syllabi/CS301.md labs P8 "sort via BST" and P12 "BST and
% traversals".)
%
% Student-facing problem statement. Instructor-only files live in
% Assignments/CS301/11-trees-bst-avl-coding/:
%   reference_solution.c  (NEVER published, gitignored)
%   test_hidden.c         (NEVER published, gitignored)
%   autograder.sh         (NEVER published, gitignored)
%   tree.h, starter_tree.c, test_visible.c  (provided to students)
%
% Fresh instances: the deck's seven-node tree (50,30,70,20,40,60,80) and the
% check-yourself tree (8,3,10,1,6,14,4,7,13) are NOT the worked example here.
% This assignment uses 9,5,12,2,7,11,15,3,6 and a second set
% 40,20,60,10,30,50,70,25,35 in the hidden suite.
%
% Compile with (macOS/Linux): cd Assignments/CS301 &&
%   TEXINPUTS=../../Preambles:$TEXINPUTS xelatex -interaction=nonstopmode 11-trees-bst-avl-coding-assignment.tex
%   BIBINPUTS=../..:$BIBINPUTS bibtex 11-trees-bst-avl-coding-assignment
%   TEXINPUTS=../../Preambles:$TEXINPUTS xelatex -interaction=nonstopmode 11-trees-bst-avl-coding-assignment.tex (x2)
% =============================================================================

\documentclass[11pt]{article}
\usepackage[margin=1in]{geometry}
% =============================================================================
% Preambles/header.tex — shared LaTeX/Beamer preamble for this project
%
% Use from a Beamer lecture by prepending:
%     \input{header}
% and compiling with TEXINPUTS=../Preambles:$TEXINPUTS (the /compile-latex
% skill does this automatically).
%
% PALETTE CONTRACT. Color names defined here MUST match the names in
% Quarto/theme-template.scss so Beamer and Quarto renderings use the same
% palette. The scripts/check-palette-sync.sh script enforces this.
%
% To customize: edit the HEX values below AND the matching $variable in
% Quarto/theme-template.scss, then run scripts/check-palette-sync.sh.
% =============================================================================

% -----------------------------------------------------------------------------
% Engine detection + encoding
%
% Our /compile-latex skill uses XeLaTeX, where inputenc is unnecessary and
% can produce encoding warnings. Only load inputenc under pdfTeX.
% -----------------------------------------------------------------------------
\usepackage{iftex}
\ifPDFTeX
  \usepackage[utf8]{inputenc}
\fi

\usepackage{amsmath, amssymb, amsthm}
\usepackage{booktabs}
\usepackage{graphicx}

% -----------------------------------------------------------------------------
% PALETTE — mirrors Quarto/theme-template.scss variable names.
% Load xcolor and define colors BEFORE anything that references them
% (notably hyperref, hypersetup, and the Beamer color assignments).
% -----------------------------------------------------------------------------
\usepackage{xcolor}

\definecolor{primary-blue}{HTML}{012169}      % headings, accents
\definecolor{primary-gold}{HTML}{B9975B}      % emphasis, borders
\definecolor{highlight-yellow}{HTML}{F2A900}  % markers, alerts
\definecolor{light-bg}{HTML}{E8EDF5}          % title / section backgrounds
\definecolor{jet}{HTML}{1A1A1A}               % body text

% Semantic colors (used in diagrams and annotations)
\definecolor{positive}{HTML}{15803D}          % good / observed / identified
\definecolor{negative}{HTML}{B91C1C}          % bad / problematic / bias
\definecolor{neutral}{HTML}{525252}           % reference / context

% Highlight accents (match .hi-* classes in the SCSS)
\definecolor{hi-slate}{HTML}{314F4F}
\definecolor{hi-green}{HTML}{2E8B57}
\definecolor{hi-red}{HTML}{C41E3A}

% Semantic aliases so skill templates and snippets can write
% \textcolor{accent}{...} without hard-coding "primary-blue".
\colorlet{accent}{primary-blue}
\colorlet{accent2}{primary-gold}

% -----------------------------------------------------------------------------
% Hyperref — loaded AFTER xcolor + palette so link colors resolve.
% -----------------------------------------------------------------------------
\usepackage{hyperref}
\hypersetup{colorlinks=true, linkcolor=primary-blue, urlcolor=primary-blue,
            citecolor=primary-blue, pdfborder={0 0 0}}

% -----------------------------------------------------------------------------
% Course code — set once per deck (after \input{header}, before \begin{document})
% via \coursecode{CS401}. Shown in the footer on every slide so a deck's course
% is identifiable without opening the title page. Empty by default.
% -----------------------------------------------------------------------------
\makeatletter
\newcommand{\coursecode}[1]{\gdef\@coursecodetext{#1}}
\gdef\@coursecodetext{}
\makeatother

% -----------------------------------------------------------------------------
% Beamer theme assignments (only apply under Beamer).
%
% We detect Beamer via \@ifclassloaded{beamer}, which is the documented,
% non-internal way. This must be wrapped in \makeatletter/\makeatother
% because \@ifclassloaded contains an @-catcode character.
% -----------------------------------------------------------------------------
\makeatletter
\@ifclassloaded{beamer}{%
  \usefonttheme{professionalfonts}
  \setbeamercolor{structure}{fg=primary-blue}
  \setbeamercolor{frametitle}{fg=primary-blue}
  \setbeamercolor{title}{fg=primary-blue}
  \setbeamercolor{subtitle}{fg=primary-gold}
  \setbeamercolor{author}{fg=primary-blue}
  \setbeamercolor{institute}{fg=neutral}
  \setbeamercolor{date}{fg=primary-gold}
  \setbeamercolor{itemize item}{fg=primary-blue}
  \setbeamercolor{itemize subitem}{fg=primary-gold}
  \setbeamercolor{alerted text}{fg=primary-gold}
  \setbeamercolor{block title}{fg=white, bg=primary-blue}
  \setbeamercolor{block body}{bg=light-bg}
  % Semantic block variants: block=definition/concept (blue), exampleblock=
  % worked example (green/positive), alertblock=key takeaway or caution (gold).
  % Keep to <=2 colored boxes per slide across ALL three variants (INV-7).
  \setbeamercolor{block title example}{fg=white, bg=positive}
  \setbeamercolor{block body example}{bg=light-bg}
  \setbeamercolor{block title alerted}{fg=white, bg=primary-gold}
  \setbeamercolor{block body alerted}{bg=light-bg}

  % Minimal footer: course code (left) + page X / N (right), no navigation clutter
  \setbeamertemplate{navigation symbols}{}
  \setbeamertemplate{footline}{%
    \color{neutral}\footnotesize\hspace{0.7em}\@coursecodetext\hfill
    \insertframenumber/\inserttotalframenumber\hspace{0.7em}\vspace{0.3em}}
}{}
\makeatother

% -----------------------------------------------------------------------------
% TikZ setup — libraries the snippet gallery uses
% -----------------------------------------------------------------------------
\usepackage{tikz}
\usetikzlibrary{arrows.meta, positioning, calc, decorations.pathreplacing,
                fit, shapes.geometric, backgrounds}

% Shared TikZ styles — snippets and hand-written diagrams can pick these up
% via `\begin{tikzpicture}[every node/.style={...}]` or by naming specific
% styles (e.g., `\node[dag-node] (X) at (0,0) {X};`).
\tikzset{
  % Node styles for causal DAGs and similar
  dag-node/.style = {
    draw=primary-blue, fill=light-bg, rounded corners=2pt,
    minimum width=1.2cm, minimum height=0.9cm, align=center, font=\small
  },
  decision-node/.style = {
    draw=primary-gold, fill=white, diamond, aspect=1.8,
    minimum width=1.8cm, minimum height=1.2cm, align=center, font=\small,
    inner sep=1pt
  },
  % Edge styles — observed vs counterfactual
  observed-edge/.style = {-{Latex[length=2.5mm]}, thick, draw=jet},
  counterfactual-edge/.style = {-{Latex[length=2.5mm]}, thick, draw=neutral, dashed},
  confound-edge/.style = {-{Latex[length=2.5mm]}, thick, draw=negative, dashed},
  % Data-point styles for plots
  observed-dot/.style = {fill=primary-blue, draw=primary-blue, circle, inner sep=1.2pt},
  counterfactual-dot/.style = {fill=white, draw=neutral, circle, inner sep=1.2pt}
}

% -----------------------------------------------------------------------------
% Convenience macros
% -----------------------------------------------------------------------------
\newcommand{\muted}[1]{\textcolor{neutral}{#1}}
\newcommand{\key}[1]{\textcolor{primary-gold}{\textbf{#1}}}
\newcommand{\good}[1]{\textcolor{positive}{#1}}
\newcommand{\bad}[1]{\textcolor{negative}{#1}}

% Standout frame macro: full-bleed dark background for section transitions.
% Beamer-only (uses \begin{frame}/\setbeamercolor/\usebeamerfont) -- do not
% call from an article-class file (e.g. Notes/<CODE>/*-notes.tex). \muted,
% \key, \good, \bad, and \coursecode above are plain xcolor/gdef and are
% safe to use from either class; verified when Notes/article-class support
% was added.
% Expand: \transitionslide{Your transition title}
\newcommand{\transitionslide}[1]{%
  \begingroup
  \setbeamercolor{background canvas}{bg=primary-blue}
  \begin{frame}[plain]
    \centering
    \vfill
    {\usebeamerfont{title}\color{white}\Large #1\par}
    \vfill
  \end{frame}
  \endgroup
}

% Section divider frame (Beamer-only), an alternative to \transitionslide with a
% darker two-line style: near-black (jet) background, a small white label line
% above a larger gold title, and a thin gold rule along the bottom edge.
% Expand: \sectiondivider{Section 2}{Pointers}
\newcommand{\sectiondivider}[2]{%
  \begingroup
  \setbeamercolor{background canvas}{bg=jet}
  \begin{frame}[plain]
    \vspace*{3.0cm}
    {\color{white}\ttfamily\large #1\par}
    \vspace{0.5cm}
    {\color{primary-gold}\LARGE #2\par}
    \begin{tikzpicture}[remember picture, overlay]
      \draw[primary-gold, line width=2.5pt]
        ($(current page.south west)+(0,0.1)$) -- ($(current page.south east)+(0,0.1)$);
    \end{tikzpicture}
  \end{frame}
  \endgroup
}

% -----------------------------------------------------------------------------
% Channel watermark — "Concepts That Click" (YouTube).
%
% Article-class files (CompetitiveExam Q/A sets, Notes, Assignments) get a
% light diagonal draft watermark on every page plus a centered footer naming
% the channel. Beamer decks get a subtle TikZ background watermark instead
% (the footline already carries the course code). Centralized here so every
% MCQ PDF is branded without per-file edits.
% -----------------------------------------------------------------------------
\makeatletter
\@ifclassloaded{article}{%
  \usepackage{draftwatermark}
  \SetWatermarkText{Concepts That Click}
  \SetWatermarkScale{1}
  \SetWatermarkFontSize{2cm}
  \SetWatermarkLightness{0.86}
  \SetWatermarkAngle{45}
  \usepackage{fancyhdr}
  \pagestyle{fancy}
  \fancyhf{}
  \fancyfoot[C]{\footnotesize\textcolor{neutral}{Concepts That Click~|~YouTube: \href{https://www.youtube.com/@concepts.thatclick}{@concepts.thatclick}}}
  \renewcommand{\headrulewidth}{0pt}
  \renewcommand{\footrulewidth}{0pt}
  \fancypagestyle{plain}{%
    \fancyhf{}%
    \fancyfoot[C]{\footnotesize\textcolor{neutral}{Concepts That Click~|~YouTube: \href{https://www.youtube.com/@concepts.thatclick}{@concepts.thatclick}}}%
    \renewcommand{\headrulewidth}{0pt}%
    \renewcommand{\footrulewidth}{0pt}%
  }
}{}
\@ifclassloaded{beamer}{%
  \setbeamertemplate{background}{%
    \begin{tikzpicture}[remember picture, overlay]
      \node[rotate=45, opacity=0.055, align=center]
        at (current page.center)
        {\Huge\textcolor{neutral}{Concepts That Click}};
    \end{tikzpicture}%
  }
}{}
\makeatother 
\coursecode{CS301}
\renewcommand{\thesection}{11.\arabic{section}}
\setlength{\emergencystretch}{6em}

\title{Coding Assignment 11: Binary Search Trees and Traversals}
\author{Auqib Hamid Lone\\[0.3em]
  \emph{Department of Computer Science \& Engineering}, \emph{GCET Ganderbal}}
\date{\today}

\begin{document}
\maketitle

\noindent\textbf{Due:} two weeks from issue. There is no automatic late penalty --- if you need more time, ask before the deadline and an extension will be granted (see the course policy in the syllabus).

\noindent This assignment autogrades Week~11's institute practicals \textbf{P8} (\key{sort via BST}) and \textbf{P12} (\key{BST and traversals}). The contract is in \texttt{tree.h} --- \emph{do not modify it}. Copy \texttt{starter\_tree.c} to \texttt{tree.c} and fill in the functions. The setting is the running calculator. In Week~10 its symbol table was a hash table: one-probe lookup, but no order. This week the same identifiers live in a binary search tree, so a sorted listing, a minimum, a maximum, and a range report all become one root-to-leaf walk.

\section{Why this problem}

A binary search tree is a binary tree with one global rule: for every node, every key in its left subtree is smaller and every key in its right subtree is larger. The rule is what turns a comparison into a decision that discards an entire subtree. Search, insert, minimum, maximum and delete are all $O(h)$ --- the cost is the \emph{height}, not the number of keys. The four traversals then expose the same tree in four different orders: inorder gives sorted output (the P8 sort), preorder copies the shape, postorder frees it safely, and level order reads it breadth-first.

The deck also shows the failure mode: inserting keys in sorted order builds a chain, and the tree degenerates to a linked list. This assignment deliberately includes that case in the hidden suite. You are implementing the plain BST, not the AVL repair; the point is to see \emph{why} the repair is needed, and to make every operation respect the ordering rule even when duplicates are present. Wirth treats tree search, insertion and deletion in Sec.~4.4.3--4.4.4, pp.135--139 \cite{Wirth2004_algorithms_data_structures}; Sedgewick's BST section starts at p.396 \cite{Sedgewick2011_algorithms}; Karumanchi's tree chapter is Ch.~6, pp.224--368 (PDF) \cite{Karumanchi2017_data_structures_algorithms_made_easy}.

\section{Your task}

Implement every function declared in \texttt{tree.h}. The tree is a \key{multiset}: the same key may occur more than once. Equal keys go in the \emph{right} subtree of an equal key, so an inorder walk prints equal keys together in nondecreasing order.

\subsection*{Creation, insertion, search}

\begin{verbatim}
Tree *tree_create(void);
void  tree_destroy(Tree *t);
int   tree_insert(Tree *t, int key);
int   tree_search(const Tree *t, int key);
size_t tree_count(const Tree *t, int key);
size_t tree_size(const Tree *t);
int   tree_height(const Tree *t);
\end{verbatim}

\texttt{tree\_insert} adds one occurrence and returns $1$ (or $-1$ on allocation failure or a NULL tree). \texttt{tree\_search} returns $1$ if the key occurs at least once, else $0$. \texttt{tree\_count} returns the number of occurrences; because equal keys are to the right, it can stop at the first equal node and count the equal run in the right subtree. \texttt{tree\_height} is measured in \emph{edges}: $-1$ for the empty tree, $0$ for one node.

\subsection*{Minimum, maximum, deletion}

\begin{verbatim}
int    tree_min(const Tree *t, int *out);
int    tree_max(const Tree *t, int *out);
int    tree_delete(Tree *t, int key);
size_t tree_remove_all(Tree *t, int key);
\end{verbatim}

\texttt{tree\_min} and \texttt{tree\_max} walk left or right to the end and return $1$ with the key in \texttt{*out} (or $0$ if the tree is empty; a NULL \texttt{out} means ``just tell me whether a key exists''). \texttt{tree\_delete} removes \emph{one} occurrence and must handle the three cases: a leaf, a node with one child, and a node with two children. The two-child case uses the \key{inorder successor}: the smallest key in the right subtree. \texttt{tree\_remove\_all} removes every occurrence and returns how many it removed.

\subsection*{Traversals, sort, range}

\begin{verbatim}
size_t tree_inorder(const Tree *t, int *out, size_t cap);
size_t tree_preorder(const Tree *t, int *out, size_t cap);
size_t tree_postorder(const Tree *t, int *out, size_t cap);
size_t tree_levelorder(const Tree *t, int *out, size_t cap);
size_t tree_sort(const int *in, size_t n, int *out, size_t cap);
size_t tree_range(const Tree *t, int lo, int hi, int *out, size_t cap);
\end{verbatim}

Each traversal writes at most \texttt{cap} keys into \texttt{out} and returns the number written. Inorder is left--node--right, preorder is node--left--right, postorder is left--right--node, and level order is breadth-first, left to right (use the Week~5 queue idea). \texttt{tree\_sort} is P8: build a BST from \texttt{in[0..n-1]} and copy its inorder into \texttt{out}; duplicates are preserved and \texttt{in} and \texttt{out} may be the same array. \texttt{tree\_range} copies every key in $[lo,hi]$, in order, and must prune subtrees that cannot contain a key in the range.

\section{Constraints}

\begin{itemize}
  \item Compile with \texttt{gcc -std=c11 -Wall -Wextra -Werror}: no errors and no warnings. No global or static mutable state.
  \item Every query is NULL-safe: a NULL tree has size $0$, height $-1$, and no keys; \texttt{tree\_insert(NULL, k)} returns $-1$; \texttt{tree\_destroy(NULL)} does nothing.
  \item Duplicates are allowed and are placed in the right subtree. \texttt{tree\_delete} removes exactly one occurrence; \texttt{tree\_remove\_all} removes all of them.
  \item Every operation must be $O(h)$ (or $O(h+k)$ for a range of $k$ keys), never $O(n)$ for a single key search on a balanced tree.
  \item The hidden suite runs $20{,}000$ keys through the tree and a $3{,}000$-key sorted insertion. The sorted case \emph{will} make the height $2{,}999$ --- that is the expected plain-BST behaviour, not a bug.
  \item The whole suite runs under \texttt{-fsanitize=address,undefined}: a leak, a double free or undefined behaviour is a failure, not a deduction.
  \item \texttt{tree\_sort} must not modify the input array. If \texttt{out} and \texttt{in} are the same buffer, build the whole tree before writing the first output key.
\end{itemize}

\section{Worked examples (these are the visible tests)}

\textbf{Example 1 --- the fresh tree.} Insert
\[
  9,\ 5,\ 12,\ 2,\ 7,\ 11,\ 15,\ 3,\ 6
\]
in that order into an empty tree. The four traversals are:
\begin{center}
{\footnotesize
\begin{tabular}{l l}
  \toprule
  inorder   & $2,3,5,6,7,9,11,12,15$ \\
  preorder  & $9,5,2,3,7,6,12,11,15$ \\
  postorder & $3,2,6,7,5,11,15,12,9$ \\
  level order & $9,5,12,2,7,11,15,3,6$ \\
  \bottomrule
\end{tabular}}
\end{center}
The height is $3$ (the path $9\to5\to2\to3$), the size is $9$, the minimum is $2$, and the maximum is $15$.

\textbf{Example 2 --- deleting a two-child node.} Delete $5$ from the tree above. Its right subtree is rooted at $7$; the leftmost node there is $6$, so $6$ replaces $5$ and the old $6$ leaf is removed. The inorder becomes
\[
  2,\ 3,\ 6,\ 7,\ 9,\ 11,\ 12,\ 15.
\]
Deleting the root $9$ next uses successor $11$ and leaves
\[
  2,\ 3,\ 6,\ 7,\ 11,\ 12,\ 15.
\]

\textbf{Example 3 --- duplicates.} Insert $4$ three times: \texttt{tree\_count(t,4)} is $3$, and \texttt{tree\_range(t,4,4)} writes three copies. Delete once: count is $2$. \texttt{tree\_remove\_all(t,4)} returns $2$ and the tree is empty.

\textbf{Example 4 --- tree sort.} With
\[
  \texttt{in} = \{4,\ -1,\ 4,\ 3,\ 0,\ -1,\ 9\},
\]
\texttt{tree\_sort} writes
\[
  -1,\ -1,\ 0,\ 3,\ 4,\ 4,\ 9.
\]
The range query $[3,9]$ writes $3,4,4,9$.

\textbf{Example 5 --- empty and NULL.} \texttt{tree\_height(NULL) = -1}, \texttt{tree\_size(NULL) = 0}, \texttt{tree\_search(NULL,1) = 0}, \texttt{tree\_create()} is not NULL, and every created tree is destroyed.

\section{How to work}

\begin{enumerate}
  \item Copy \texttt{starter\_tree.c} to \texttt{tree.c}; keep \texttt{tree.h} unchanged.
  \item Build and run the visible tests:
\begin{verbatim}
gcc -std=c11 -Wall -Wextra -Werror -o test_visible test_visible.c tree.c
./test_visible
\end{verbatim}
  \item Get Examples 1 and 3 green before writing deletion. A wrong traversal order makes every later comparison look like a delete bug.
  \item Trace the two-child deletion on paper: find the successor, copy its key, then delete the successor from the right subtree. The successor has no left child, so the recursive call is Case 1 or 2.
  \item Run the same build under the sanitizers before submitting:
\begin{verbatim}
gcc -std=c11 -Wall -Wextra -Werror -g \
    -fsanitize=address,undefined -o test_visible test_visible.c tree.c
./test_visible
\end{verbatim}
  \item The hidden suite adds: randomized insert/delete/search/count against an array model, forced duplicate runs, all three deletion cases, a sorted $3{,}000$-key chain, a $20{,}000$-key scale run, range boundaries, cap truncation, in-place tree sort, and leak checking.
\end{enumerate}

\section{Submission and grading}

Submit \texttt{tree.c} through the course submission form. Grading is by the autograder --- correctness on the hidden suite, compiled with \texttt{-Werror} and run under \texttt{-fsanitize=address,undefined} --- plus a short code read for four structural requirements: a real BST (not an array in disguise), all three deletion cases with a correct successor, four genuine traversals (level order uses a queue, not recursion), and a tree sort that preserves duplicates. A solution that is correct on the small cases but never frees nodes is partially credited; a warning under \texttt{-Werror}, a leak, or a sanitizer report is a failure.

\bibliography{../../Bibliography_base}
\bibliographystyle{plain}
\end{document}
